
\section{Introduction}

\begin{frame}{maillage}

\begin{columns}
\begin{column}{0.45 \textwidth}

\begin{block}{Loi de conservation scalaire}
\begin{equation}
\left\{
\begin{aligned}
    &\partial_t u + \partial_x f(u) = 0\qquad \forall (x,t) \in \mathbb{D} \times \R^+ \\
    & u : \mathbb{D} \times \R^+ \to \Omega \subset R^M
\end{aligned}    
\right.
\end{equation}
\end{block}
\end{column}

\begin{column}{0.5 \textwidth}
\begin{align*}
    f(u) &= c\,u            & \text{advection}\\
    f(u) &= \frac{4u^2}{4u^2 +(1-u)^2}  &\text{Buckley}\\
    f(u) &= (\rho v,\,\, \rho v^2 + p, \,\,\rho e\, v+ p\, v ) &\text{Euler}
\end{align*}
\end{column}

\end{columns}

\begin{columns}
\begin{column}{0.45\textwidth}
\begin{block}{critères d'un maillage de $\mathbb{D}$}
\begin{itemize}
    \item $\mathbb{D} = \bigcup_i \omega_i$ \qquad $\omega_i \neq \emptyset $\\
    \item $\omega_i \cap \omega_j = \emptyset$ si $i\neq j$\\
    \item $\omega_i$ est convexe 
\end{itemize}
\end{block}
\end{column}

\begin{column}{0.45\textwidth}
\begin{block}{maillage 1D}
\begin{itemize}
    \item $\omega_i = [x_{i-1/2},x_{i+1/2}]$ \\
    \item si $\mathbb{D} = [a,b]$, $x_{1/2}=a$ et $x_{N_x+1/2} =b$
\end{itemize}
\end{block}
\end{column}
\end{columns}

\begin{block}{Condition de Courant–Friedrichs–Lewy }
Condition régissant le maillage espace-temps dépendant de la solution et de sa vitesse de déplacement dans le maillage espace
\begin{tabular}{p{6cm}c}
\begin{equation}
    \Delta t \leq \frac{\Delta x}{2 \gamma}\qquad \gamma = \max_u |f'(u)|
\end{equation}
&
\raisebox{-15mm}{\usebox{\Riemann}}
\end{tabular}   
\end{block}

\end{frame}

% ---------------------------------------------------------
% ---------------------------------------------------------

\begin{frame}{Solveurs de Riemann approchée}

\begin{block}{Solveur de Riemann Approchée}
On approche ${u_r}(\frac{x}{t},u_L,u_R)$ solution d'un porblème de Riemann par $\widetilde{u}_r$.\\
le SRA sera consitent au sens de la moyenne si : $\Delta t,\, \Delta x$ respectent la condition CFL et :
\begin{tabular}{p{2.5cm}| p{7cm}}
\begin{equation}
    \widetilde{u}_r(\frac{x}{t},u,u)=u 
\end{equation}
&
\hspace{1cm}
\begin{equation}
    \frac{1}{\Delta x}\int_{-\Delta x/2}^{\Delta x/2} \widetilde{u}_r(\frac{x}{\Delta t},u_L,u_R) dx
    = \frac{u_L + u_R}{2} -\frac{\Delta t}{2\Delta x}\left(f(u_R)-f(u_L)\right)
\end{equation}
\end{tabular}
\end{block}

\begin{columns}

\begin{column}{0.4 \textwidth}
\begin{block}{Solveur de Rusanov/Lax}
\begin{equation}
    \left\{
    \begin{aligned}
        &\mathcal{F}(u,v) = \frac{f(u)+f(v)}{2} - \frac{\gamma}{2}(v-u) \\
        &u^* = \frac{u_L +u_R}{2} -\frac{1}{2\gamma}(f(u_R)-f(u_L))
    \end{aligned}
    \right.
\end{equation}
\end{block}
\end{column}

\begin{column}{0.55 \textwidth}
\begin{block}{Solveur de Godunov}
\begin{equation}
    \left\{
    \begin{aligned}
        &\mathcal{F}(u,v) = \frac{f(u)+f(v)}{2} - \frac{\gamma_L}{2}(u^*-u_L) -\frac{\gamma_R}{2}(u_R-u^*)  \\
        &u^* = \frac{\gamma_R\,u_R-\gamma_L\,u_L }{\gamma_R - \gamma_L} -\frac{1}{\gamma_R - \gamma_L}(f(u_R)-f(u_L))
    \end{aligned}
    \right.
\end{equation}
\end{block}
\end{column}

\end{columns}

\end{frame}
% ==========================================================================
% ==========================================================================
\section{Schéma Volumes Finis}

\begin{frame}{Schéma Volumes Finis}

\begin{block}{SRA entropique}
On dit que le SRA est entropique si il respecte cette condition supplémentaire.
\begin{tabular}{p{2cm}| p{7cm}}
\begin{align*}
\partial_t S(u) + \partial_x G(u) \leq 0 \\ \forall (x,t) \in \R \times \R^+
\end{align*}
&
\begin{align*}
    \frac{1}{\Delta x}\int_{-\Delta x/2}^{\Delta x/2} S(\widetilde{u}_r(\frac{x}{\Delta t},u_L,u_R)) dx 
   = \frac{S(u_L) + S(u_R)}{2}-  \frac{\Delta t}{2\Delta x}\left(G(u_R)-G(u_L)\right)
\end{align*}
\end{tabular}
\end{block}
\begin{columns}
\begin{column}{0.35 \textwidth}
\begin{block}{Schéma semi-discrétisé}
\begin{align}    
    & \partial_t \overline{u_j} = \frac{-1}{|\omega_j|} (\mathcal{F}_{j+1/2}-\mathcal{F}_{j-1/2})\\
    & \overline{u}_j^n = \int_{\omega_j} u(x,t^n)dx \\
    &\mathcal{F}_{j+1/2} = \mathcal{F}(u_j,u_{j+1})
\end{align}
\end{block}
\end{column}

\begin{column}{0.6 \textwidth}
\begin{block}{Schéma discrétisé par Euler Explicite}
\begin{equation}  
    \overline{u_j}^{n+1} = \overline{u_j}^n - \frac{\Delta_t}{|\omega_j|} (\mathcal{F}_{j+1/2}-\mathcal{F}_{j-1/2})
\end{equation}
\end{block}

\begin{block}{Schéma RK-SSP}
% \begin{align*}  
    $u^1 = \hspace{1.5mm}u^0 + \hspace{5mm}\Delta t L(u^0)$ \\
    $u^2 = \alpha u^0 +(1-\alpha)(u^1 + \Delta t L(u^1))$\hspace{.2cm} ordre 2 : $\alpha = \frac{1}{2}$ 
    $u^3 = \beta u^0 +(1-\beta)(u^2 + \Delta t L(u^2))$\hspace{.2cm} ordre 3 : $\alpha = \frac{3}{4}$, $\beta = \frac{1}{3}$
% \end{align*}
\end{block}

\end{column}
\end{columns}



\end{frame}

% ---------------------------------------------------------
% ---------------------------------------------------------

\begin{frame}{Problème : dissipation de la solution}


\red{insert results: }\\
\red{left : sinus advection, sinus burgers (bon comportement)}\\
\red{right: rectangulaire advection, composite advection(dissipation)}

\end{frame}
